% discussion.tex (v6): six paragraphs, no subheadings (Sci Rep), about 850 words.
\section{Discussion}\label{sec:discussion}

The cell-year design answers the four questions with one set of inputs, splits and scores. Year-specific forcing improves a physical model where observations span many years; the validation split decides whether direct learning looks best or fails; a residual learner on a Stefan anchor gains two to three centimetres in two of three regions and nothing in the third; and on the same cells and years the combined model is at least as accurate as the published products once training-site leakage is removed. None of the individual components is new. The contribution is the time definition, the common test bed and the conditions under which each component helps.

The time definition is the part most often left implicit. Existing statistical products use multi-year mean climate and multi-year mean site values, so a map of theirs has no year and their validation statistics mix the interannual and the spatial parts of the error. Accumulating degree-days to the observation day matters because the observations themselves are not end-of-season values: within one cell-year, September measurements are about 10 cm deeper than August measurements, and a model trained on full-year forcing attributes this to the site rather than to the date. The ratio of observation-day to full-year degree-days (0.73 in Alaska) shows how much of the season the conventional input overstates. The same reasoning explains why the Alaska year effect cannot be assessed: with half of the labels in 2014, the yearly and climatological forcings are nearly collinear with the site effect, and no split based on these data can separate them. Designing campaigns so that sites are revisited in several years would resolve this more than any modelling choice.

The split results repeat, with a yearly design, what has been shown for spatial prediction in general\cite{Roberts2017,Ploton2020,MeyerPebesma2022}: random splits credit a learner with information that it will not have at a new place. Here the optimism is \noptimLeDz{}--\noptimCaDz{} cm, and under region hold-out direct learning is no better than, and in the Lena Delta far worse than, the Stefan coefficient transferred from other regions. Product validations reported on random site splits are therefore not comparable with ours and are likely optimistic. The physics anchor is what makes the combined model safe: its region-hold-out error stays within a few centimetres of the physical model, and its gain inside a region is bounded by the shrinkage factor on the residual.

The conditions of the gain are consistent across regions. The residual learner needs terrain and soil covariates; with climate variables only it is worse than the physical model in Alaska, which indicates that the residual it learns from climate alone is site noise. The gain is larger in the warmer terciles of each region, where a single regional Stefan coefficient cannot follow the spread of thaw depth between organic-rich, shaded and well-drained cells, and small in the cold terciles, where thaw is shallow and its spatial variation is within the measurement noise. Two negative results are as informative as the positive ones. Physics pseudo-labels did not improve on direct learning in any region, because they carry no information beyond the physical model that already enters the residual design. Changing the learner moved the gain by less than the differences between regions and never changed its sign, so the choice of learner is secondary to the choice of inputs and of the physical anchor.

The product comparison shows why same-cell, same-year scoring with a leakage mask is necessary. Products trained on the Alaska observations reach low errors on the cells they were trained on and higher errors elsewhere; once cells within 25 km of their training sites are removed, the ordering changes. In the Lena Delta, where no product was trained, the combined model is better than every product by 10 cm or more. We note that the products were not designed for yearly values and that part of their error on single years is interannual variation that they do not attempt to represent.

The main limitation is the interannual skill. Correlations between predicted and observed year-to-year change stay below 0.3 for our methods in Alaska and the Lena Delta, the residual learner adds no interannual skill to the physical anchor, and at the Lena Delta sites with repeated observations the Wei et al. product follows year-to-year change better than any of our methods. Matching the forcing to the observation day removes part of the year-wise bias in Alaska, but the largest remaining yearly biases are explained by changes in campaign composition rather than by climate. Reanalysis degree-days at 9 km resolution, the absence of snow and soil-moisture observations at the site scale, and the dominance of single-visit sites all contribute. The maps therefore represent the year of their forcing in their regional pattern but should not be read as site-level interannual series outside the observed range. Two directions follow. Longer site records and campaigns that revisit cells would make the year effect identifiable in Alaska; and the residual design can take yearly satellite inputs (snow-off date, surface deformation from radar interferometry) that we tested only as static summaries. The year-labelled maps with their area of applicability are intended as a baseline against which such additions can be scored on the same cells and years.
